\documentclass[11pt]{article}
\usepackage{mathblatt}
% Präambel des Setzers (werkzeuge/setzer-praeambel.tex, wortgleich); dazu \gteil (Muster T2U)
\makeatletter
\setlength{\mbpfbe}{0mm}% keine Punkte (pruefung.md § 5: nur beim Prüfstein)
\newlength{\gEin}\setlength{\gEin}{0pt}
\let\gSchrift\relax
\renewcommand{\pfkreuzzeile}[1]{\par\noindent\hangindent3.2em\hangafter1\hspace*{1.6em}$\square$\ #1\par}
\newcommand{\gkreuz}[1]{$\square$\ #1\hspace{2em}}
\newcommand{\gkreuzzeile}[2]{\par\noindent\hangindent2.2em\hangafter1\mbox{#1\,$\square$}\ #2\par}% Buchstabe, Option: eine je Zeile (neben Skizze)
% links, was man ansieht, rechts, was man tut; Spaltengrenze überall an derselben Stelle (gemeinsam.md „Satz“)
\newsavebox{\gbox}\newlength{\gLinks}
% v0.10 (Lehrer 13.10.): keine starre Halbierung – die linke Spalte ist so breit wie die Figur (+ 6 mm), die
% Antwort steht direkt daneben; ist die Figur breiter als 0,62 Zeile, steht die Antwort darunter
\newcommand{\gzwei}[2]{\par\smallskip\noindent\sbox{\gbox}{#1}%
  \setlength{\gLinks}{\dimexpr\wd\gbox+6mm\relax}%
  \ifdim\gLinks>0.62\linewidth
    \usebox{\gbox}\par\smallskip\noindent\begin{minipage}{\linewidth}#2\end{minipage}%
  \else
    \begin{minipage}[t]{\gLinks}\vspace{0pt}\usebox{\gbox}\end{minipage}%
    \begin{minipage}[t]{\dimexpr\linewidth-\gLinks\relax}\vspace{0pt}\raggedright #2\end{minipage}%
  \fi\par}
% v0.9 (12.10., pruefung.md § 5 „Marke mit B, Bedingungen rechts“): drittes Argument von pfaufg = Bedingungen
% der Prüfung (v0.10: \pfohnetr = TR durchgestrichen, „FOR“), klein grau im rechten Rand auf der Grundlinie der Nummer, Zeilen mit \\
\newcommand{\gbedrand}[2]{\ifblank{#2}{}{\rlap{\hspace*{\dimexpr#1+2mm\relax}%
  \parbox[t]{15mm}{\raggedright\scriptsize\leavevmode\color{mbgrau}\strut#2}}}}
% Teil eines Bündels (a, b …): Marke links im Rand, Bedingungen rechts im Rand, auf der Zeile des Teils
\newcommand{\gteilrand}[2]{\llap{\mbpfrandmarke{#1}\hspace*{\mbpfnr}}\gbedrand{\linewidth}{#2}}
% Teil eines Bündels: Buchstabe in eigener Spalte, Inhalt eingerückt (\gEin = Einzug, damit die Spaltengrenze von
% \gzwei an derselben Stelle bleibt)
\newenvironment{gbteil}[3]{\par\smallskip\noindent\gteilrand{#2}{#3}\makebox[\mbpfteil][l]{\textbf{#1}}%
  \begin{minipage}[t]{\dimexpr\linewidth-\mbpfteil\relax}\setlength{\gEin}{\mbpfteil}\raggedright\strut\ignorespaces}%
 {\end{minipage}\par}
% Grafik höchstens so breit wie die Spalte (Bündel nebeneinander, Koordinatensystem ohne Karo)
\newsavebox{\gpbox}
\newcommand{\gpassend}[1]{\sbox{\gpbox}{#1}\ifdim\wd\gpbox>\linewidth\typeout{GPASSEND: verkleinert \the\wd\gpbox\space auf \the\linewidth}\resizebox{\linewidth}{!}{\usebox{\gpbox}}\else\usebox{\gpbox}\fi}
% Skizzenfeld (v0.10, Lehrer 13.10.): Karo ab dem linken Rand über die Zeile, „Skizze“ klein grau oben links im Feld
% (wie \pfskizzenfeld der Vorlage); der rechte Rand bleibt den Bedingungen
\newcommand{\gskizze}[1]{\par\smallskip\noindent\begin{tikzpicture}
  \clip (0,0) rectangle ({\linewidth-1mm},{#1*10mm});
  \draw[black!16,line width=0.3pt,step=5mm] (0,0) grid ({\linewidth},{#1*10mm});
  \node[anchor=north west,inner sep=2pt,font=\scriptsize,text=mbgrau,fill=white] at (0,{#1*10mm}) {Skizze};
  \end{tikzpicture}\par}
\RenewDocumentEnvironment{pfaufg}{m m m}{%
  \begin{lrbox}{\mbaufgabebox}\begin{minipage}[t]{\linewidth}%
  \gSchrift\noindent\hspace*{\gEin}\mbpfrandmarke{#1}\gbedrand{\linewidth-\gEin}{#3}%
  \begin{minipage}[t]{\mbpfnr}\strut\textbf{#2}\end{minipage}%
  \begin{minipage}[t]{\dimexpr\linewidth-\gEin-\mbpfnr\relax}\raggedright\strut\ignorespaces}%
 {\end{minipage}%
  \end{minipage}\end{lrbox}%
  \setlength{\mbaufgabehoehe}{\dimexpr\ht\mbaufgabebox+\dp\mbaufgabebox\relax}%
  \ifdefempty{\gkopfnach}{\par\addvspace{5pt}\Needspace*{\mbaufgabehoehe}}{%
    % Stufenkopf, der mit dieser (ersten) Aufgabe zusammenbleibt: Platz für beide verlangen, dann den Kopf setzen
    % (12.10., Quadratische: ein Kopf stand allein unten, die Aufgabe mit Koordinatensystem auf der nächsten Seite)
    \par\Needspace*{\dimexpr\mbaufgabehoehe+40pt\relax}\gkopfnach\gdef\gkopfnach{}\par\addvspace{5pt}}%
  \noindent\usebox{\mbaufgabebox}\par\addvspace{5pt}}
% Stufenkopf vormerken: \gkopfvor{<Kopf>} – die nächste pfaufg setzt ihn, sobald beide zusammen Platz haben
\def\gkopfnach{}
\newcommand{\gkopfvor}[1]{\gdef\gkopfnach{#1}}
% Teilaufgabe mit Skizze links in gzwei: Buchstabe in eigener Spalte, darunter Text und Skizze
\newcommand{\gteil}[2]{\makebox[\mbpfteil][l]{\textbf{#1}}\parbox[t]{\dimexpr0.5\textwidth-\gEin-\mbpfnr-5mm-\mbpfteil\relax}{\raggedright #2}}
% Stufenkopf der Vorlage (mathblatt.sty 2026-10-12: Linie, Abstand, 13 pt, \pfjahre) ohne ihren Vorhalt von
% 14 Zeilen: der Setzer fordert vor dem Kopf 0,16 Seitenhöhe (Kopf + erste Aufgabe), mehr nicht
\setlength{\mbpfstufevorhalt}{0pt}
\makeatother
\newcommand{\karo}[1]{\par\smallskip\noindent\begin{tikzpicture}
  \clip (0,0) rectangle ({\linewidth-1mm},{#1*10mm});
  \draw[black!16,line width=0.3pt,step=5mm] (0,0) grid ({\linewidth},{#1*10mm});
  \end{tikzpicture}\par}
\newcommand{\kk}[1]{$\square$\,#1\hspace{0.9em}}
\newcommand{\linie}[1]{\,{\color{black!45}\rule[-1pt]{#1}{0.4pt}}\,}
\newcommand{\pkt}{\textperiodcentered{}}
% v0.7 (12.10.): Schülerfrage unter dem Stufenkopf; Ankreuzoption als Bild (Buchstabe und Kästchen darunter)
\newcommand{\gfrage}[1]{\par\nobreak\noindent\hspace*{\mbpfnr}{\small\itshape #1}\par\nobreak\addvspace{2pt}}
\newlength{\gbildbreite}\setlength{\gbildbreite}{0.32\linewidth}% Setzer setzt sie je Reihe (vier Bilder: 0,24)
\newcommand{\gbildoption}[3]{\begin{minipage}[t][#3][b]{\gbildbreite}\centering #2\par\smallskip\mbox{#1\,$\square$}\end{minipage}}
% Lineare (12.10.): Bild mit Buchstabe ohne Kästchen (Diagramm zuordnen); Aussage mit „wahr / falsch“ zum Ankreuzen
\newcommand{\gbildname}[3]{\begin{minipage}[t][#3][b]{\gbildbreite}\centering #2\par\smallskip\mbox{\textbf{#1}}\end{minipage}}
\newcommand{\gwahrfalsch}[2]{\par\noindent\hangindent2.2em\hangafter1\makebox[2.2em][l]{(#1)}#2\hfill\mbox{wahr $\square$\quad falsch $\square$}\par}
\begin{document}

\pfheftstil
\pflaufkopf[P10 \pkt{} Prüfung]{Wachstum}{74A2}\pfseitevon
\pfheftkopf{Wachstum}{}
\pfstufenzeile{Tabelle ergänzen \pkt{} Graph zeichnen \pkt{} Gleichung gesucht \pkt{} Graph zuordnen}

% ===================== Stufe: Tabelle ergänzen =====================
\Needspace*{0.16\textheight}\pfstufekopf{Tabelle ergänzen}{\pfjahre{26}{20 \pkt{} 19 \pkt{} 17 \pkt{} 16}}
% Schülerfrage der Stufe (Gliederung; Einstieg nur von dort, nie aus der Bank – 12.10.)
\begin{pfaufg}{}{1.}{}
Ein Betrag wächst jedes Jahr um denselben Prozentsatz. Mit welcher Zahl multiplizierst du ihn, um den Betrag im nächsten Jahr zu bekommen?
\par\smallskip
\par\noindent\makebox[\mbpfteil][l]{\textbf{a)}}Wachstum um 5 \%\hfill\linie{30mm}
\par\noindent\makebox[\mbpfteil][l]{\textbf{b)}}Wachstum um 2,3 \%\hfill\linie{30mm}
\par\fusshilfe{\mbox{1:~a) 1,05; b) 1,023}}
\end{pfaufg}
% 2026-FOR-K7a-m1 (Original 2026-FOR-K7a, Schwere I, rang 1; herausgelöste Fassung: keine)
\begin{pfaufg}{nach~P10\\’26~\pkt{}~7a}{2.}{}
Can bezieht Anfang 2026 eine neue Wohnung. Laut Mietvertrag steigt die Miete jedes Jahr um 2,3 \%. Im ersten Jahr zahlt er 590 € im Monat. Can legt eine Tabelle an. Vervollständige sie.
\par\smallskip\noindent\begin{tabular}{|l|c|c|c|c|}\hline Jahr & 2026 & 2027 & 2028 & 2029\\\hline Miete in € & \leerzelle & 603,57 & 617,45 & \leerzelle\\\hline\end{tabular}\par
\karo{2}
\fusshilfe{\mbox{2:~2026: 590,00 €; 2029: 631,65 €}}
\end{pfaufg}

% ===================== Stufe: Punkte darstellen =====================
\gkopfvor{\pfstufekopf{Graph zeichnen}{\pfjahre{25}{21 \pkt{} 17 \pkt{} 16}}}
% kein Einstieg: keine Schülerfrage in der Gliederung (Schülerfrage, -Fälle, -Form, -Weg, -Skizze; oder Frage) – Einstieg nur von dort, nie aus der Bank (pruefung.md § 5, 12.10.)
% 2025-OS-K7a-m1 (Original 2025-OS-K7a, Schwere I, rang 2; herausgelöste Fassung: keine)
\begin{pfaufg}{nach~P10\\’25~\pkt{}~7a}{3.}{}
Man beobachtet einen Bakterienstamm, der exponentiell wächst. Die Tabelle zeigt die Zahl der Bakterien in den ersten 4 Stunden. Stelle das Wachstum der Bakterien im Koordinatensystem dar.
\par\smallskip\noindent\begin{minipage}[t]{0.55\linewidth}\vspace{0pt}\begin{tabular}[t]{|l|c|c|c|c|c|}\hline Zeit in Stunden & 0 & 1 & 2 & 3 & 4\\\hline Anzahl der Bakterien & 70 & 91 & 118 & 154 & 200\\\hline\end{tabular}\end{minipage}\hfill\begin{minipage}[t]{0.4\linewidth}\vspace{0pt}\begin{ksys}[xmin=0,xmax=5,ymin=0,ymax=250,xstep=1,ystep=25,karo=0.42,xlabel=Zeit in h,ylabel=Anzahl]\end{ksys}\end{minipage}\par
\fusshilfe{\mbox{3:~fünf Punkte (0|70), (1|91), (2|118), (3|154), (4|200), zunehmend steiler ansteigend}}
\end{pfaufg}

% ===================== Stufe: Faktor bestimmen, Gleichung aufstellen =====================
\Needspace*{0.16\textheight}\pfstufekopf{Gleichung gesucht}{\pfjahre{26 \pkt{} 25}{20 \pkt{} 19 \pkt{} 18 \pkt{} 17 \pkt{} 16}}
% Schülerfrage der Stufe (Gliederung; Einstieg nur von dort, nie aus der Bank – 12.10.)
\begin{pfaufg}{}{4.}{}
Eine Größe wächst in einem Jahr von 200 auf 210. Mit welcher Zahl wurde sie multipliziert, und um wie viel Prozent ist sie gewachsen?
\par\smallskip
\par\noindent von 200 auf 210\hfill\hfill Faktor: \linie{30mm}\par\noindent\hfill Wachstum: \linie{30mm}\,\%
\par\fusshilfe{\mbox{4:~210 : 200 = 1,05; 5 \%}}
\end{pfaufg}
% 2026-FOR-K7c-m1 (Original 2026-FOR-K7c, Schwere II, rang 3; herausgelöste Fassung: keine)
\begin{pfaufg}{nach~P10\\’26~\pkt{}~7c}{5.}{FOR}
Can bezieht Anfang 2026 eine neue Wohnung. Laut Mietvertrag steigt die Miete jedes Jahr um 2,3 \%. Im ersten Jahr zahlt er 590 € im Monat. Stelle eine Gleichung auf, mit der man die Monatsmiete nach beliebig vielen Jahren berechnen kann. Bestimme die Monatsmiete im Jahr 2040.
\karo{3}
\fusshilfe{\mbox{5:~590 · $1{,}023^{t}$; 811,17 €}}
\end{pfaufg}
% 2025-OS-K7b-m1 (Original 2025-OS-K7b, Schwere II, rang 4; herausgelöste Fassung: keine)
\begin{pfaufg}{nach~P10\\’25~\pkt{}~7b}{6.}{FOR}
Man beobachtet einen Bakterienstamm, der exponentiell wächst. Die Tabelle zeigt die Zahl der Bakterien in den ersten 4 Stunden. Zeige, dass der Wachstumsfaktor 1,3 ist. Gib eine Gleichung an, die die Zahl der Bakterien nach $t$ Stunden beschreibt. Prüfe mit einer Rechnung, ob diese Aussage stimmt: „Nach 24 Stunden sind es mehr als 35000 Bakterien.“
\par\smallskip\noindent\begin{tabular}[t]{|l|c|c|c|c|c|}\hline Zeit in Stunden & 0 & 1 & 2 & 3 & 4\\\hline Anzahl der Bakterien & 70 & 91 & 118 & 154 & 200\\\hline\end{tabular}\par
\karo{4}
\fusshilfe{\mbox{6:~1,3); 70 · $1{,}3^{t}$; 37996}}
\end{pfaufg}

% ===================== Stufe: Graph zuordnen und begründen =====================
\gkopfvor{\pfstufekopf{Graph zuordnen}{\pfjahre{26}{19}}}
% kein Einstieg: keine Schülerfrage in der Gliederung (Schülerfrage, -Fälle, -Form, -Weg, -Skizze; oder Frage) – Einstieg nur von dort, nie aus der Bank (pruefung.md § 5, 12.10.)
% 2026-FOR-K7b-m1 (Original 2026-FOR-K7b, Schwere III, rang 5; herausgelöste Fassung: keine)
\begin{pfaufg}{nach~P10\\’26~\pkt{}~7b}{7.}{}
Can bezieht Anfang 2026 eine neue Wohnung. Laut Mietvertrag steigt die Miete jedes Jahr um 2,3 \%. Im ersten Jahr zahlt er 590 € im Monat. Entscheide, welcher Graph zur Entwicklung der Miete passt, und kreuze ihn an. Erkläre bei den zwei übrigen Graphen jeweils, weshalb sie nicht passen.
\par\smallskip\noindent\setlength{\gbildbreite}{\dimexpr0.96\linewidth/3\relax}\sbox{\gbox}{\vbox{\begin{tikzpicture}[x=0.5cm,y=0.5cm,line width=0.5pt]\draw[black!20,line width=0.3pt,step=0.5] (0,0) grid (4.5,3.5);\draw[-{Stealth[length=1.5mm]}] (0,0)--(4.8,0);\draw[-{Stealth[length=1.5mm]}] (0,0)--(0,3.8);\node[below,font=\tiny] at (2.4,0) {Zeit in Jahren};\node[right,font=\tiny] at (0,3.8) {Miete in €};\draw[line width=0.8pt,domain=0:4.2,samples=40,smooth] plot (\x,{0.3*(exp(0.55*\x)-1)});\end{tikzpicture}\par\begin{tikzpicture}[x=0.5cm,y=0.5cm,line width=0.5pt]\draw[black!20,line width=0.3pt,step=0.5] (0,0) grid (4.5,3.5);\draw[-{Stealth[length=1.5mm]}] (0,0)--(4.8,0);\draw[-{Stealth[length=1.5mm]}] (0,0)--(0,3.8);\node[below,font=\tiny] at (2.4,0) {Zeit in Jahren};\node[right,font=\tiny] at (0,3.8) {Miete in €};\draw[line width=0.8pt,domain=0:4.2,samples=40,smooth] plot (\x,{0.9*exp(0.3*\x)});\end{tikzpicture}\par\begin{tikzpicture}[x=0.5cm,y=0.5cm,line width=0.5pt]\draw[black!20,line width=0.3pt,step=0.5] (0,0) grid (4.5,3.5);\draw[-{Stealth[length=1.5mm]}] (0,0)--(4.8,0);\draw[-{Stealth[length=1.5mm]}] (0,0)--(0,3.8);\node[below,font=\tiny] at (2.4,0) {Zeit in Jahren};\node[right,font=\tiny] at (0,3.8) {Miete in €};\draw[line width=0.8pt] (0,0.9)--(4.2,3.1);\end{tikzpicture}}}\gbildoption{A}{\begin{tikzpicture}[x=0.5cm,y=0.5cm,line width=0.5pt]\draw[black!20,line width=0.3pt,step=0.5] (0,0) grid (4.5,3.5);\draw[-{Stealth[length=1.5mm]}] (0,0)--(4.8,0);\draw[-{Stealth[length=1.5mm]}] (0,0)--(0,3.8);\node[below,font=\tiny] at (2.4,0) {Zeit in Jahren};\node[right,font=\tiny] at (0,3.8) {Miete in €};\draw[line width=0.8pt,domain=0:4.2,samples=40,smooth] plot (\x,{0.3*(exp(0.55*\x)-1)});\end{tikzpicture}}{\dimexpr\ht\gbox/3+\baselineskip\relax}\hfill\gbildoption{B}{\begin{tikzpicture}[x=0.5cm,y=0.5cm,line width=0.5pt]\draw[black!20,line width=0.3pt,step=0.5] (0,0) grid (4.5,3.5);\draw[-{Stealth[length=1.5mm]}] (0,0)--(4.8,0);\draw[-{Stealth[length=1.5mm]}] (0,0)--(0,3.8);\node[below,font=\tiny] at (2.4,0) {Zeit in Jahren};\node[right,font=\tiny] at (0,3.8) {Miete in €};\draw[line width=0.8pt,domain=0:4.2,samples=40,smooth] plot (\x,{0.9*exp(0.3*\x)});\end{tikzpicture}}{\dimexpr\ht\gbox/3+\baselineskip\relax}\hfill\gbildoption{C}{\begin{tikzpicture}[x=0.5cm,y=0.5cm,line width=0.5pt]\draw[black!20,line width=0.3pt,step=0.5] (0,0) grid (4.5,3.5);\draw[-{Stealth[length=1.5mm]}] (0,0)--(4.8,0);\draw[-{Stealth[length=1.5mm]}] (0,0)--(0,3.8);\node[below,font=\tiny] at (2.4,0) {Zeit in Jahren};\node[right,font=\tiny] at (0,3.8) {Miete in €};\draw[line width=0.8pt] (0,0.9)--(4.2,3.1);\end{tikzpicture}}{\dimexpr\ht\gbox/3+\baselineskip\relax}\par
\karo{2}
\fusshilfe{\mbox{7:~B; jedes Jahr gleicher Betrag statt gleicher Prozentsatz}}
\end{pfaufg}

\pfblattende{11}{26 FOR & 25 & 24 & 23 & 22 & 21 & 20 & 19 & 18 & 17 & 16}{7a 7b 7c & 7a 7b &  &  &  & 6b & 4a 4c & 7a 7b 7c & 2a & 7a 7b 7c 7d & 4a 4b 4e}%
{\pfausschnitt{}{Wachstum}{}}
\end{document}